\section*{Chapitre \ref{ch:b1:structure-spectroscopy} --- La spectroscopie au service de la structure}

\begin{solution}{exo:b1:structure-spectroscopy:1}
$\varepsilon = A/(\ell c) = 0.64/(1.00 \times \num{2.0e-5}) =
\qty{3.2e4}{L/(mol.cm)}$. Concentration triplée~: $A = 1.92$, si la loi
reste valable à cette \omterm{def:b1:structure-spectroscopy:absorbance}{absorbance}.
\end{solution}

\begin{solution}{exo:b1:structure-spectroscopy:2}
\ce{C6H6}~: $D = (12 + 2 - 6)/2 = 4$, le benzène (un cycle et trois
liaisons $\pi$). \ce{C4H8O2}~: 1, l'éthanoate d'éthyle ou l'acide
butanoïque. \ce{C3H7NO}~: $(6 + 2 + 1
- 7)/2 = 1$, le propanamide. \ce{C2H3Cl}~: 1, le chloroéthène.
\ce{C10H14O}~:
$(20 + 2 - 14)/2 = 4$~: la carvone possède un cycle, deux C=C et un C=O.
\end{solution}

\begin{solution}{exo:b1:structure-spectroscopy:3}
Propanone~: un signal (6H). Acide éthanoïque~: deux (3H, 1H).
2-Méthylpropane~: deux (9H, 1H). 1,2-Dichloroéthane~: un (4H).
\end{solution}

\begin{solution}{exo:b1:structure-spectroscopy:4}
$\delta = 1460/400 = \qty{3.65}{ppm}$~; à \qty{90}{MHz}, le signal se
trouve à $3.65 \times 90 = \qty{329}{Hz}$ de la référence. Les
\omterm{def:b1:structure-spectroscopy:coupling}{constantes de couplage} restent les mêmes en hertz tandis que les écarts
de \omterm{def:b1:structure-spectroscopy:chemical-shift}{déplacement chimique} croissent avec le champ~: des \omterm{def:b1:structure-spectroscopy:coupling}{multiplets} qui se
chevauchent à bas champ se séparent à haut champ.
\end{solution}

\begin{solution}{exo:b1:structure-spectroscopy:5}
\ce{CH3}~: triplet (3H)~; \ce{CH2} central~: cinq voisins, sextuplet
1\,:\,5\,:\,10\,:\,10\,:\,5\,:\,1 (2H)~; \ce{CH2Br}~: triplet (2H), le
plus déblindé.
\end{solution}

\begin{solution}{exo:b1:structure-spectroscopy:6}
Éthanol~: O--H (3666) et pas de C=O~; propanone~: C=O (1738) et pas de
O--H~; acide éthanoïque~: les deux, O--H (3582) et C=O (1788).
\end{solution}

\begin{solution}{exo:b1:structure-spectroscopy:7}
Éthanoate d'éthyle~: quadruplet (2H) vers 4.1, singulet (3H) vers 2.0,
triplet (3H) vers 1.3 ppm. Propanoate de méthyle~: singulet (3H) de
\ce{OCH3}, déblindé par l'oxygène (vers 3.7 ppm), quadruplet (2H) de
\ce{CH2C=O} (vers 2.3), triplet (3H) vers 1.1. Le quadruplet est le
signal le plus déblindé dans le premier et pas dans le second~: un
quadruplet vers 4 ppm signale \ce{O-CH2CH3}.
\end{solution}

\begin{solution}{exo:b1:structure-spectroscopy:8}
Raies distantes de \qty{0.080}{ppm}, soit $0.080 \times 90 = \qty{7.2}{Hz}$~:
la \omterm{def:b1:structure-spectroscopy:coupling}{constante de couplage}. Le triplet présente le même écart~: le
\ce{CH2} et le \ce{CH3} sont couplés l'un à l'autre.
\end{solution}

\begin{solution}{exo:b1:structure-spectroscopy:9}
Une bande C=O et un seul signal~: la propanone, \ce{CH3COCH3}. L'isomère
aldéhyde, le propanal \ce{CH3CH2CHO}, présente la bande d'élongation C--H
du groupe \ce{CHO} et son proton très déblindé.
\end{solution}

\begin{solution}{exo:b1:structure-spectroscopy:10}
Le premier ensemble scinde la raie en $n_1 + 1$ raies~; chacune d'elles
est scindée par le second ensemble en $n_2 + 1$~: $(n_1 + 1)(n_2 + 1)$
raies. Si $J_1 =
J_2$, la position d'une raie ne dépend que de $k_1 + k_2$, qui prend $n_1 + n_2 +
1$ valeurs~: c'est la règle des $n + 1$ avec $n = n_1 + n_2$. \ce{CH2}
central du 1-bromopropane~: $4 \times 3 = 12$ raies en général, six (un
sextuplet) lorsque les deux constantes sont égales.
\end{solution}

\begin{solution}{exo:b1:structure-spectroscopy:11}
$D = 0$~; un alcool (bande O--H, pas de C=O). Doublet (6H)~: deux
\ce{CH3} équivalents portés par un même CH~; \omterm{def:b1:structure-spectroscopy:coupling}{multiplet} (1H)~: ce CH~;
doublet (2H)~: un \ce{CH2} voisin du CH et portant le OH~; singulet
élargi~: OH. Le 2-méthylpropan-1-ol, \ce{(CH3)2CH-CH2-OH}.
\end{solution}

\begin{solution}{exo:b1:structure-spectroscopy:12}
Chaque espèce absorbe indépendamment~: $A = A_1 + A_2 = \ell(\varepsilon_1c_1
+ \varepsilon_2c_2)$. À deux longueurs d'onde, les quatre coefficients
étant connus, deux équations linéaires donnent $c_1$ et $c_2$.
\end{solution}

\begin{solution}{pb:b1:structure-spectroscopy:1}
\textbf{1.} $0.5690 \times 12.0/44.0 = \qty{0.1552}{g}$.
\textbf{2.} $0.2328 \times 2.0/18.0 = \qty{0.02587}{g}$.
\textbf{3.} $0.2500 - 0.1552 - 0.0259 = \qty{0.0689}{g}$.
\textbf{4.} C \qty{62.1}{\percent}, H \qty{10.3}{\percent}, O
\qty{27.6}{\percent}.
\textbf{5.} C $0.1552/12.0 = 0.01293$, H $0.02587$, O $0.0689/16.0 =
0.00431$ mol~: rapport $3.00 : 6.00 : 1$, \ce{C3H6O}.
\textbf{6.} $M = \rho RT/p = 3740 \times 8.314 \times 373.15/(1.000 \times 10^5) =
\qty{116}{g/mol}$ (masse volumique en \unit{g/m^3}).
\textbf{7.} $116/58 = 2$~: \ce{C6H12O2}, $D = (12 + 2 - 12)/2 = 1$.
\textbf{8.} Une élongation C=O.
\textbf{9.} Toute liaison O--H~: ni acide, ni alcool.
\textbf{10.} Une élongation C--O, intense~: le C--O d'un ester.
\textbf{11.} Élongations C--H de groupes \ce{CH2} et \ce{CH3}.
\textbf{12.} Une seule liaison $\pi$, celle du C=O~; avec un C--O intense
et pas de O--H, un ester. Un acide exige un O--H~; un diol n'a pas de
C=O. Une cétone portant un groupe éther aurait son C=O près des
\qty{1738}{cm^{-1}} de la propanone, alors que les \qty{1754}{cm^{-1}}
observés sont proches des \qty{1764}{cm^{-1}} de l'éthanoate d'éthyle~;
la RMN de la partie III confirme l'ester.
\textbf{13.} Cinq~; $2 + 2 + 2 + 3 + 3 = 12$, les douze hydrogènes.
\textbf{14.} Un \ce{CH2} à trois voisins (un \ce{CH3}), lié à l'oxygène
(déblindé)~: \ce{-O-CH2-CH3}.
\textbf{15.} Le \ce{CH3} de ce groupe éthyle, à deux voisins~; les deux
signaux sont couplés l'un à l'autre, d'où un même $J$.
\textbf{16.} Un \ce{CH2} à deux voisins, voisin du C=O.
\textbf{17.} Un \ce{CH2} à cinq voisins~: entre un \ce{CH2} et un
\ce{CH3}.
\textbf{18.} Un \ce{CH3} à deux voisins, loin du groupe caractéristique.
\textbf{19.} \ce{CH3CH2O-} et \ce{-CH2CH2CH3}.
\textbf{20.} Les protons à 4.13 ppm sont portés par le carbone lié à
l'oxygène, ceux à 2.28 ppm sont voisins du C=O~: deux groupes
électroattracteurs.
\textbf{21.} \ce{CH3CH2-O-CO-CH2CH2CH3} ou \ce{CH3CH2-CO-O-CH2CH2CH3}.
Seule la première place sur l'oxygène le \ce{CH2} à trois voisins (le
quadruplet à 4.13 ppm).
\textbf{22.} Le propanoate de propyle est la seconde~: son \ce{OCH2}
serait un triplet et le quadruplet se trouverait vers 2.3 ppm.
\textbf{23.} L'éthanoate de butyle, \ce{CH3CO-O-C4H9}, présenterait un
singulet (3H) vers 2.0 ppm et pas de quadruplet.
\textbf{24.} Le butanoate d'éthyle, \ce{CH3CH2CH2-CO-O-CH2CH3}.
\textbf{25.} IR~: C=O 1754, C--O 1182, C--H 2982, pas de O--H. RMN~:
\ce{OCH2} 4.13 (q), \ce{CH2CO} 2.28 (t), \ce{CH2} central 1.66
(sextuplet), \ce{CH3} de l'ester 1.25 (t), \ce{CH3} de la chaîne
0.95 (t), \omterm{def:b1:structure-spectroscopy:integration}{intégrations} 2, 2, 2, 3, 3.
\textbf{26.} \textbf{\qty{116}{g/mol}}~: le butanoate d'éthyle,
\ce{C6H12O2}.
\end{solution}
